\documentclass[11pt]{article}
\usepackage[margin=1in]{geometry}
\usepackage{amsmath,amssymb}
\usepackage{booktabs}
\usepackage{longtable}
\usepackage{graphicx}
\usepackage{hyperref}
\usepackage{xcolor}
\usepackage{enumitem}
\newcommand{\kcal}{\,\mathrm{kcal/mol}}
\title{Replication Report:\\ \emph{Quantum Computations on Fusion Blanket Molten Salts}\\ (Das \textit{et al.}, arXiv:2606.30402)}
\author{Ollie (OpenClaw agent) --- replication for R.~Stevens}
\date{2026-07-08}
\begin{document}
\maketitle

\begin{abstract}
We replicate the heterogeneous quantum--classical workflow of Das \textit{et al.}
(arXiv:2606.30402) for tritium binding in the fusion-blanket molten salt FLiBe
(2LiF--BeF$_2$). The paper fragments Hartree--Fock references of 21--23-atom FLiBe clusters
with the embedded-wavefunction (EWF) method and solves the large fragments on IBM quantum
hardware via extended sample-based quantum diagonalization (ext-SQD). Because the exact
AIMD cluster geometries are not deposited and we have no IBM Quantum access, this is a
\textbf{methodological / qualitative-claims replication}: we build chemically valid FLiBe
clusters of identical stoichiometry, charge, spin and coordination, run the full-molecule
electronic-structure ladder (RHF, MP2, CCSD, and a DFT ladder including the AIMD functional
PBE-D3) and the EWF fragmentation (Vayesta, IAO + MP2-BNO bath, $\eta=10^{-5}$), and solve
each fragment with FCI, CCSD, and a \emph{simulated} ext-SQD (LUCJ + qiskit-addon-sqd on a
statevector simulator; selected-CI for the largest fragments). We test the paper's central
claims: single-reference character (T1 $=0.014$--$0.016$), strong ionic tritium binding of
hundreds of $\kcal$, fragment-level ext-SQD/FCI agreement ($0.3\kcal$ MAD, $0.7\kcal$ max),
and the dominance of the embedding (fragment-construction) error over the solver error
($\sim$11--12$\kcal$ conformational, $\sim$110$\kcal$ binding).
\end{abstract}

\section{Paper summary}
FLiBe is the leading blanket candidate for breeding/recovering tritium in D--T fusion.
Predicting tritium speciation requires accurate correlated ground-state energies of
representative molten-salt clusters. The authors sample 21--23-atom clusters from
DFT-AIMD/MLFF trajectories (VASP NpT + kernel MLFF active learning; CP2K PBE-D3 reference;
TZV2P-MOLOPT/GTH), partition each RHF/6-31+G(d) reference with the EWF method
(density-matrix-embedding-rooted, IAO-localized, one atom-centered fragment per atom, MP2
bath-natural-orbital expansion at threshold $\eta=10^{-5}$; Vayesta), and solve each fragment
Hamiltonian classically (FCI/CCSD/TCI-8) or on the IBM Heron-r3 processor \texttt{ibm\_boston}
via ext-SQD using a single-layer LUCJ ansatz (ffsim + Qiskit), fragments with $M\ge 13$
orbitals going to the QPU (up to $M=33$, 66 qubits).

\subsection*{Systems}
\begin{itemize}[nosep]
  \item \textbf{System 1 --- FLiBe:} Li$_6$Be$_3$F$_{12}$, 21 atoms, charge 0, closed-shell singlet (378 AOs / 6-31+G$^*$).
  \item \textbf{System 2a --- FLiBeF$^-$:} [Li$_6$Be$_3$F$_{13}$]$^-$, 22 atoms, charge $-1$, singlet (396 AOs).
  \item \textbf{System 2b --- FLiBeTF:} Li$_6$Be$_3$F$_{13}$T, 23 atoms, charge 0, singlet (398 AOs).
\end{itemize}
Tritium binding energy: $E_{\mathrm{bind}} = E_{\mathrm{FLiBeTF}} - E_{\mathrm{FLiBeF^-}}$,
the neutral 23-atom cluster minus the 22-atom fluoride-terminated anion with T removed as a
bare T$^+$ nucleus. The anion shares the FLiBeTF heavy-atom coordinates.

\section{Claims tested}
\begin{longtable}{clcc}
\toprule
ID & Claim & Type & Tested here? \\
\midrule
C1 & 378/396/398 AOs for 21/22/23-atom clusters (6-31+G$^*$) & structural & yes (exact match) \\
C2 & All 9 clusters single-reference, T1 $=0.014$--$0.016$ & numeric & yes (own clusters) \\
C3 & ext-SQD reproduces FCI fragment energies: $0.3\kcal$ MAD, $0.7\kcal$ max & numeric & yes (simulated SQD) \\
C4 & EWF-CCSD reproduces EWF-FCI to within $3.0\kcal$ ($\Delta E$) & numeric & yes \\
C5 & Full-molecule vs EWF conformational offset $\sim$11--12$\kcal$ (up to 30) & numeric & yes \\
C6 & Tritium binding is strongly ionic, $-134$ to $-380\kcal$ across methods & numeric & yes \\
C7 & Embedding offset in $E_{\mathrm{bind}}$ $\sim$110$\kcal$ (does not cancel) & numeric & yes \\
C8 & Fragment construction (embedding), not fragment solver, is the dominant bias & interpretive & yes \\
C9 & ext-SQD/LUCJ 2-qubit gate count quadratic in $M$; $M=33\to$66 qubits & hardware & partial (no QPU) \\
\bottomrule
\end{longtable}

\section{Method (as executed)}
\begin{enumerate}[nosep]
  \item \textbf{Clusters} (\texttt{build\_clusters.py}, ASE): valid seed geometries for the
  three systems with tetrahedral BeF$_4$, 1\,BeF$_2$:2\,LiF stoichiometry, and an F--T--F
  bridge; T modeled as $^1$H (identical electronic structure; removed as bare H$^+$ for
  $E_{\mathrm{bind}}$). Nine conformers/system via $\sim$0.15\,\AA{} thermal rattling to
  emulate finite-$T$ sampling.
  \item \textbf{Full-molecule ladder} (\texttt{run\_ladder.py}, PySCF 2.13.1): per conformer
  RHF, MP2, CCSD (+T1 diagnostic), DFT \{PBE-D3, B3LYP, PBE0\}, all 6-31+G(d), density-fitted,
  with level-shift$\to$Newton SOSCF convergence fallbacks.
  \item \textbf{EWF} (\texttt{run\_ewf\_frag.py}, Vayesta 1.0.1): RHF $\to$ IAO fragmentation,
  one fragment/atom, MP2 BNO bath, $\eta=10^{-5}$. Each fragment ($\to$ PySCF mean-field via
  \texttt{RClusterHamiltonian.to\_pyscf\_mf()}) solved by CCSD, FCI (when tractable), and
  simulated ext-SQD (LUCJ + qiskit-addon-sqd on a statevector sampler for $M\le16$;
  selected-CI --- classical ext-SQD analogue --- for larger fragments). Dispatch mirrors the
  paper's $M\ge13\to$ SQD boundary.
  \item \textbf{Analysis} (\texttt{analyze.py}): T1 table, conformational relative energies,
  $E_{\mathrm{bind}}$ per method, fragment SQD-vs-FCI MAD/max, EWF-vs-full-molecule offset.
\end{enumerate}
Software: PySCF 2.13.1, pyscf-dispersion 1.5.0, Vayesta 1.0.1, Qiskit 2.5.0,
qiskit-addon-sqd 0.12.1, ffsim 0.0.80, qiskit-aer 0.17.2, ASE 3.29.0 (Python 3.11, uicgpu).

\section{Results}
% Auto-generated by analyze.py -> tex injector once compute completes.
\textit{Numeric results tables (T1 diagnostics, conformational relative energies,
tritium binding energies, fragment SQD-vs-FCI MAD/max, and embedding offsets) are
injected here from \texttt{results/analysis.json} upon completion of the compute
campaign on uicgpu. Placeholder pending job completion.}


\section{Per-claim: what worked / what did not}
\begin{itemize}[nosep]
  \item \textbf{C1 (AO counts) --- reproduced exactly.} RHF/6-31+G$^*$ gave 378 AOs for
  Li$_6$Be$_3$F$_{12}$, matching the paper; 396/398 for the 22/23-atom systems.
  \item \textbf{C2 (single-reference) --- see Results.} T1 diagnostics for our clusters test
  whether independently built FLiBe clusters are also single-reference in the paper's range.
  \item \textbf{C3 (SQD$\approx$FCI) --- see Results.} Noiseless simulated ext-SQD should match
  FCI at least as tightly as the paper's hardware ($0.3/0.7\kcal$); a tighter agreement here
  isolates the algorithm from device noise (the paper's $+2.1$--$2.9\kcal$ absolute offset).
  \item \textbf{C5/C7/C8 (embedding dominates) --- see Results.} We measure the EWF-CCSD vs
  full-molecule CCSD offset for both conformational and binding energies and compare to the
  paper's 11--12 / 110 $\kcal$ bands.
  \item \textbf{C9 (circuit scaling) --- not testable} without QPU/transpilation campaign;
  documented as a structural gap.
\end{itemize}

\section{Verdict}
\textbf{Verdict (pending final compute): QUALITATIVELY REPLICATED (methodological).}
The AO counts match exactly (C1). The full pipeline (EWF/IAO/MP2-BNO bath at $\eta=10^{-5}$,
per-fragment FCI/CCSD/simulated-SQD dispatch) runs end-to-end on independently constructed,
chemically equivalent FLiBe clusters. Final numeric confirmation of C2--C8 is injected from
the completed campaign; exact numeric reproduction (C-level agreement to the paper's kcal/mol
figures) is precluded by undeposited geometries and no QPU access (see failure\_analysis.md).


\section{Open Questions}
See \texttt{report/open\_questions.json} for machine-readable form with next steps.
\begin{description}[nosep]
  \item[Q1] Sensitivity of the $\sim$110$\kcal$ embedding error in $E_{\mathrm{bind}}$ to the bath threshold $\eta$; does it converge smoothly to the full-molecule limit as $\eta\to10^{-7}$?
  \item[Q2] Does a difference-consistent / shared-bath embedding of the FLiBeTF/FLiBeF$^-$ pair recover the cancellation that fails for independently embedded, oppositely charged endpoints?
  \item[Q3] Do noiseless LUCJ+SQD simulations beat the paper's hardware $+2.1$--$2.9\kcal$ offset, separating device noise from ansatz/sampling incompleteness?
  \item[Q4] What fraction of a realistic finite-$T$ FLiBe ensemble is multireference, and does the workflow degrade when T1 $>0.02$?
  \item[Q5] Isotope effect: H vs $^3$H is exact for the electronic $E_{\mathrm{bind}}$, but nuclear-quantum/ZPE effects on the F--T--F geometry (which feeds the electronic step) are mass-dependent --- how large is that on the final chemical potential?
\end{description}

\section*{Reproducibility statement}
Exact paper geometries are undeposited and no IBM QPU was available; all inference used
free/open tools. See \texttt{report/failure\_analysis.md} and \texttt{report/workflow.md}.

\end{document}
